为在题给数据下得到可闭合、可复现的计算模型，本文采用以下假设。它们同时构成模型的适用范围。

\begin{enumerate}
  \item 药材可视为均质、各向同性的直圆柱，输出对应长度中点横截面；忽略轴向端部换热与传质。
  \item 第一至三问几何固定，$R=0.02\,\mathrm{m}$、$L=0.25\,\mathrm{m}$；第四问长度仍取定值，仅半径按给定轨迹作均匀比例收缩。
  \item 初始温度与干基含水率在药材内均匀，分别为 $28^\circ\mathrm{C}$ 与 $2.55\,\mathrm{kg/kg}$。
  \item 烘房温度与水分浓度由附件观测点经 PCHIP 插值进入边界，精确通过全部观测点，不对曲线作平滑或删点。
  \item 表面换热、传质采用牛顿冷却型第三类边界，系数取题给常数；空气含湿量与药材含水率之差作为等效传质驱动力，不另引入未给出的吸附等温线。
  \item 主计算模型不显式计入蒸发潜热、辐射和体积热源。表面失水计入质量平衡，但汽化耗热不计入能量方程。
  \item 水分扩散采用题给经验公式。第一问 $D=D(C)$；第二、三问 $D=D(C,T)$；第四问更换另一组 $D(C,T)$。温度进入指数时使用局部开尔文温度。
  \item 附件给出的 $\rho(C)$ 在能量方程中作为有效热密度使用，不把它同时解释为可变干物质密度。
  \item 第三、四问以计算域内最大含水率首次低于 $0.15\,\mathrm{kg/kg}$ 作为烘干结束判据；4\,h 后烘房取 $50^\circ\mathrm{C}$、$0.05\,\mathrm{kg/kg}$ 作为恒温延续，该延续不是未来实测。
\end{enumerate}
